\section{The full level record on the relative cover}

% BEGIN L73D1
\paragraph{L73D1. Lifting a framed point with its section generator}
Use the forgetful map $r:X_2\to X$ of the relative higher model at
the levels of RTG10. Write $m$ for the cartesian map of universal
polarized abelian schemes, and use the upper and lower coefficient
records with their specified common field $M$. Given an algebraically
closed field $k$ with the two compatible $M$-algebra structures and
an original framed point $y$ of $X$, the fiber of $r$ over its point
$x$ is a nonempty finite \'etale $k$-scheme. Choose $x_2$ in this
fiber. Its structure map is the required one for the upper
coefficient record, since the two embeddings of $M$ agree.

The cartesian square identifies the two abelian varieties over $k$
by an isomorphism $p:A_{x_2}\to A_x$ with the specified projections.
Pull each tangent vector of $y$ back by $p^{-1}$, using dual numbers
for tangent vectors. This preserves the zero section, addition,
scalar multiplication and the endomorphisms from $\mathcal O_E$.
The eigenvalue in the upper field maps to the same element of $k$
as the lower eigenvalue because the coefficient comparison is the
identity on $M$. The basis test for each eigenspace transports in both
directions. Thus the resulting ordered vectors form the original type
of eigenframe.

The cartesian polarization supplies the isomorphism $w$ from the
upper theta line to the pullback of the lower one, with its prescribed
generator at the identity. Choose the upper section $s_2$ by transporting
$p^*s$ through $w^{-1}$, where $s$ is the actual section of $y$.
It generates because both maps are isomorphisms. In the notation
of \texttt{FrOver}, take its intermediate section $t$ to be the
pullback of $s$ in the pulled-back lower trivializations. Then
$w(s_2)=t$ gives \texttt{SecMaps} on every overlap, and in each
lower chart $i$ the local representative of $t$ is the pullback
of $s_i$ along $p$, restricted to the indicated upper open. This is
the last chartwise equality in \texttt{FrOver}. The same $w$
respects the rigidifying generators by the cubical pullback law.
Together with $r\circ x_2=x$ and the
pushforwards of the eigenvectors, these are all the fields of the
required framed lift.
% END L73D1

% BEGIN L73D2
\paragraph{L73D2. A level datum with the canonical connection and unit section}
For the given lower model and every tuple level of RTG10, the relative
higher model yields an original \texttt{DeltaDat} with its prescribed
integer $d$, ideal $I$, upper level, model and coefficient record,
satisfying \texttt{IsLev} with character section $1$. Choose the
star-ring structure and the canonical connection on this upper model
by CNC15. For the provisional finite function of \texttt{DeltaDat}
use the constant $1$; \texttt{IsLev} places no condition on it.
RTG10 supplies the divisibility, integer-level product, coprimality,
half-ideal product and sum-with-conjugate conditions. By construction
$[d]\eta_2(a)=\eta(a)$ after forgetting the upper level. Since both
auxiliary levels equal one, their ratio is one and both associated
auxiliary marks are trivial. This proves both \texttt{EtaOK}
identities. L73D1 supplies the algebraically closed framed-point
clause.

The constant rule $1$ is natural for every test morphism and its
triangular character at weight zero and index zero is one. To verify
its nonzero evaluation recall that \texttt{evR} returns zero outside
the locus of points descending from a test ring. Every upper framed
point over an algebraically closed $k$ is in that locus. Here TDR13
uses base change for the first theta power on the given family.
Indeed the theta line on each fiber is ample and principally
polarizing, so its space of sections has dimension one by
Riemann--Roch and its higher cohomology vanishes in characteristic
zero. The proper smooth derived pushforward is perfect and commutes
with base change for this line, which is flat over the base.
Its fiber cohomology in degree zero only makes it a line bundle in
degree zero locally, by splitting off the invertible differentials
of a finite free complex. Thus its sections commute with arbitrary
base change, including the maps used in framed-point descent.
TDR13 now gives an original test presentation of the given point.
In that presentation $1$ maps to $1\in k$. This completes
\texttt{IsLev}, including its last nonvanishing clause. T73K3 can
replace this unit section with the prescribed character section when
the marked upper dictionary and component-field hypotheses are
available.
% END L73D2
